\documentclass{beamer}

\newcommand{\week}{Intro 2 of 2}

\title{Overview of the Financial System}
\subtitle{Financial Instruments, Financial Intermediaries\\
Reference: Mishkin, Chapter 2}
\author{Econ 357}
\date{\week}

% Reference the shared preamble
\input{preamble}

% Set the footer -- change 
\setbeamertemplate{footline}{
  \leavevmode%
  \vspace{2ex}
  \hbox{%
    % Left box: Econ 357
    \begin{beamercolorbox}[wd=.4\paperwidth,ht=2.5ex,dp=1ex,left]{author in head/foot}%
      \hspace{1em}Econ 357
    \end{beamercolorbox}%
    % Middle box: Week
    \begin{beamercolorbox}[wd=.2\paperwidth,ht=2.5ex,dp=1ex,center]{date in head/foot}%
      \centering\week
    \end{beamercolorbox}%
    % Right box: Slide numbers
    \begin{beamercolorbox}[wd=.4\paperwidth,ht=2.5ex,dp=1ex,center]{date in head/foot}%
      \hfill\insertframenumber{} 
    \end{beamercolorbox}%
  }%
  \vskip0pt%
}

\begin{document}

\frame{\titlepage}

\begin{frame}
    \frametitle{Outline}
    \begin{enumerate}
        \item Financial Instruments
        \begin{itemize}
            \item Treasury Debt, Corporate Debt, Equities
            \item Primary v. Secondary Market
        \end{itemize}
        \item Financial Intermediaries
        \begin{itemize}
            \item Functions
            \item Examples
        \end{itemize}
        \item Practice Questions
    \end{enumerate}
\end{frame}

\begin{frame}[t]
    \frametitle{1. Financial Instruments}
    \section{Financial Instruments}
    \framesubtitle{Treasury Debt}

    Issuer (borrower):
    \vspace{3em}

    Purchaser (lender): 
    \vspace{3em}

    Types of Instruments:

\end{frame}

\begin{frame}[t]
    \frametitle{1. Financial Instruments}
    \framesubtitle{Corporate Debt}

    Issuer (borrower):
    \vspace{3em}

    Purchaser (lender): 
    \vspace{3em}

    Types of Instruments:

\end{frame}

\begin{frame}[t]
    \frametitle{1. Financial Instruments}
    \framesubtitle{Public Equities}

    Issuer (borrower):
    \vspace{3em}

    Purchaser (lender): 
    \vspace{3em}

    Types of Instruments:

\end{frame}

\begin{frame}[t]
    \frametitle{1. Financial Instruments}
    \section{Primary v. Secondary Markets}
    \framesubtitle{Primary v Secondary Markets}

        \begin{center}
    \textbf{Primary and Secondary Markets}
    \end{center}
    \begin{tblr}{
      colspec = { Q[l,wd=3cm] Q[c,wd=2cm] |Q[c,wd=2cm] |Q[c,wd=2cm] }, vline{3} = {1-Z}{}
    }
        \toprule
        & US Treasuries & Corp Bonds & US Equities \\
        \midrule
        Outstanding (as of 2024) & \$28 tr & \$11.1 tr  & \$62 tr  \\
        Issuance (as of 2024) & \$4.6 tr & \$1.9 tr  & $\sim$\$200 bn  \\
        Trading (ave daily vol 2023) & $\sim$\$900 bn & \$46 bn & \$516 bn \\
        \bottomrule
    \end{tblr}
    \vspace{1em}

    \blfootnote{Source: SIFMA, SEC, FINRA}

\end{frame}

\begin{frame}[t]
    \frametitle{1. Financial Instruments}
    \framesubtitle{Primary v Secondary Markets}

    There is substantially more fixed income issuance than equity issuance every year.
    \vspace{0.5em}

    \centering
    \includegraphics[width=0.8\textwidth]{Econ357/figures/01_intro_2_issuance.png}

    \blfootnote{Source: SIFMA}

\end{frame}

\begin{frame}[t]
    \frametitle{1. Financial Instruments}
        \framesubtitle{Primary v Secondary Markets}

    The majority of fixed income issuance is US Treasuries
    \vspace{0.5em}

    \centering
    \includegraphics[width=0.7\textwidth]{Econ357/figures/01_intro_2_fixed_issuance.png}

    \blfootnote{Source: SIFMA}

\end{frame}

\begin{frame}[t]
    \frametitle{1. Financial Instruments}
        \framesubtitle{Primary v Secondary Markets}

    Equity issuance (IPO = "initial public offering") is rare and tends to cluster.
    \vspace{0.5em}

    \centering
    \includegraphics[width=0.7\textwidth]{Econ357/figures/01_intro_2_ipos.png}

    \blfootnote{Source: SEC}

\end{frame}

\begin{frame}[t]
    \frametitle{1. Financial Instruments}
        \framesubtitle{Primary v Secondary Markets}

    Corporate bond issuance is large and regular
    \vspace{0.5em}

    \centering
    \includegraphics[width=0.7\textwidth]{Econ357/figures/01_intro_2_corp_issuance.png}

    \blfootnote{Source: SEC}

\end{frame}

\begin{frame}[t]
    \frametitle{1. Financial Instruments}
        \framesubtitle{Primary v Secondary Markets}

    Equities are heavily traded
    \vspace{0.5em}

    \centering
    \includegraphics[width=0.7\textwidth]{Econ357/figures/01_intro_2_trading.png}

    \blfootnote{Source: FINRA}

\end{frame}

\begin{frame}[t]
    \frametitle{1. Financial Instruments}
        \framesubtitle{Primary v Secondary Markets}

    Equitiy trading has increased over time
    \vspace{0.5em}

    \centering
    \includegraphics[width=0.7\textwidth]{Econ357/figures/01_intro_2_trading4.png}

    \blfootnote{Source: IMF}

\end{frame}

\begin{frame}[t]
    \frametitle{1. Financial Instruments}
        \framesubtitle{Primary v Secondary Markets}

    ...but has been a relatively constant share of GDP
    \vspace{0.5em}

    \centering
    \includegraphics[width=0.7\textwidth]{Econ357/figures/01_intro_2_trading5.png}

    \blfootnote{Source: IMF}

\end{frame}

\begin{frame}[t]
    \frametitle{2. Financial Intermediaries}
    \section{Financial Intermediaries}

        \centering
    \includegraphics[width=0.9\textwidth]{Econ357/figures/01_intro_2_intermediaries.jpg}

    \blfootnote{Source: Mishkin}

\end{frame}

\begin{frame}[t]
    \frametitle{2. Financial Intermediaries}

    Reasons financial intermediaries exist:
    \begin{itemize}
    \item Transaction Costs
    \item Risk Transformation
    \item Information Costs
    \end{itemize}

\end{frame}

\begin{frame}[t]
    \frametitle{2. Financial Intermediaries}   

    \textbf{Transaction Costs}
    \vspace{1em}

    Definition:
    \vspace{1em}
    
    Examples in financial transactions:
    \vspace{1em}

    Discussion:
    \begin{itemize}
        \item How do financial intermediaries reduce transaction costs?
        \item Economies of scale
        \item Examples of financial intermediaries that reduce transaction costs.
    \end{itemize}

\end{frame}

\begin{frame}[t]
    \frametitle{2. Financial Intermediaries}

    \textbf{Risk Transformation}
    \vspace{1em}

    Definition:
    \vspace{2em}
    
    Examples in financial transactions:
    \vspace{2em}

    Discussion:
    \begin{itemize}
        \item How do financial intermediaries transform risk?
        \item Risk sharing, asset transformation, maturity transformation, diversification
        \item Examples of financial intermediaries that transform risk.
    \end{itemize}

\end{frame}

\begin{frame}[t]
    \frametitle{2. Financial Intermediaries}

    \textbf{Information Costs}
    \vspace{1em}

    Definition:
    \vspace{2em}
    
    Examples in financial transactions:
    \vspace{2em}

    Discussion:
    \begin{itemize}
        \item Especially important when there is asymmetric information, which can cause problems with \textbf{moral hazard} and \textbf{adverse selection}
        \item Intermediaries can address asymmetric information through underwriting and monitoring
        \item Examples of financial intermediaries that transform risk.
    \end{itemize}

\end{frame}

\begin{frame}[t]
    \frametitle{2. Financial Intermediaries}
    \framesubtitle{Examples}

    \vspace{2em}
    
        \centering
    \includegraphics[width=0.9\textwidth]{Econ357/figures/01_intro_2_intermediaries2.png}

\end{frame}

\begin{frame}[t]
    \frametitle{2. Financial Intermediaries}
    \framesubtitle{Examples}
    
    \textbf{Types of Financial Intermediaries:}
    \begin{itemize}
        \item Depository Institutions: Commercial Banks, Credit Unions
        \item Contractual Savings Institutions: Insurance Companies, Pension Funds
        \item Investment Intermediaries: Mutual Funds, ETFs, Hedge Funds, Investment Banks, Brokerage
    \end{itemize}
    \vspace{1em}
    
    How does each type address the problems of transaction costs, information costs, or risk?

\end{frame}

\begin{frame}[t]
    \frametitle{2. Financial Intermediaries}
    \framesubtitle{Examples}
    
    Three Specific Examples (perhaps relevant to you…):

    \begin{enumerate}
        \item Checking account (e.g Wells Fargo, Bank of California)
        \item Brokerage account (e.g. Robinhood, Charles Schwab)
        \item Retirement Account or Life Insurance (e.g. Vanguard or Lincoln Financial)
    \end{enumerate}
    \vspace{1em}

    What intermediaries are involved?   What types of financial instruments are they buying?

\end{frame}

\begin{practiceframe}[t]
    \frametitle{3. Practice}
    \section{Practice}

    \textbf{Review}
    \begin{itemize}
        \item Define real asset and financial asset
        \item Financial assets are always a liability AND an asset
        \item Balance sheet: write down ”T-charts” for households, corporations, and US government
        \item Basic characteristics of stocks and bonds: who is the issuer?  Buyer?  Market size and trading frequency
        \item Why are financial intermediaries necessary?
        \item Identify financial intermediaries that address the following issues: risk, transaction costs, information costs, asymmetric information (moral hazard and adverse selection)
    \end{itemize}

\end{practiceframe}

\begin{practiceframe}[t]
    \frametitle{3. Practice}

    \begin{enumerate}
        \item Who issues Treasury Bonds?   Equities?  Corporate Bonds?
        \item What is the largest financial liability for the household sector?
        \item What is an example of a real asset on a corporate balance sheet?
        \item A subsidiary is a company owned by another company.  Is the subsidiary a real or a financial asset for the parent company?
    \end{enumerate}


\end{practiceframe}

\begin{practiceframe}[t]
    \frametitle{3. Practice}


Which of the following statements about the characteristics of debt and equity is FALSE?
\begin{enumerate}[A)]
    \item They can both be long-term financial instruments.
    \item They can both be short-term financial instruments.
    \item They both involve a claim on the issuer's income.
    \item They both enable a corporation to raise funds.
\end{enumerate}
\vspace{1em}

Equity holders are a corporation's (BLANK). That means the corporation must pay all of its debt holders before it pays its 
\begin{enumerate}[A)]
   \item equity holders.
    \item debtors
    \item brokers
    \item residual claimants
    \item underwriters
\end{enumerate}

\end{practiceframe}

\begin{practiceframe}[t]
    \frametitle{3. Practice}

    \begin{enumerate}
    \setcounter{enumi}{4}
        \item How does adverse selection explain why you are more likely to make a loan to a family member than a stranger?
        \item If you are an employer, what types of moral hazard do you worry about with regards to your employees?
        \item What problem is being addressed by a mutual fund that selects stocks and then pursues a “buy-and-hold” strategy?   How would you describe its strategy?
        \item Define “maturity transformation”
    \end{enumerate}

\end{practiceframe}

\end{document}